\chapter{تنظيم نظام التشغيل}
\label{CH:FIRST}

تتمثل إحدى المتطلبات الأساسية والتطبيقة في تصميم أنظمة التشغيل في تحقيق
\indextext{العزل القوي} (Strong Isolation).
يقتضي العزل منع العمليات من إتلاف أو التجسس أو التعديل على ذاكرة وحالات العمليات الأخرى،
وكذلك حماية نظام التشغيل نفسه من البرامج الخبيثة أو المعطوبة التي تعمل في فضاء المستخدم.
في هذا الفصل، نبين كيف ينظم \texttt{xv6} نظام التشغيل حول هذا الهدف،
والطريقة التي تعتمد بها النواة على العتاد المادي لتحقيق العزل، وكيفية إقلاع النواة وتفريغ العملية الأولى.

\section{تجريد الموارد الفيزيائية}
\label{sec:abstracting_resources}

منهجًا إبتدائيًا للتفاعل مع الموارد المادية هو السماح لبرامج التطبيقات بالوصول المباشر إلى العتاد الصلب.
على سبيل المثال، يمكن لكل تطبيق استخدام تعليمات الذاكرة المباشرة أو الوصول لبايتات القرص الصلب.
ولكن هذا النهج المباشر يفشل في تقديم العزل؛
إذ يمكن لبرنامج معطوب أن يكتب فوق ذاكرة برنامج آخر أو يسيطر على القرص الصلب للوصول لمحتويات البرامج الأخرى.
وحتى بدون أخطاء، فإن البرامج التي تتشارك العتاد المادي المباشر تجعل التطبيقات صعبة البناء للغاية.

لذلك، يحل نظام التشغيل هذه المشاكل عن طريق حظر الوصول المباشر للعتاد،
وبدلاً من ذلك يقدم نظام التشغيل \textbf{تجريدات محكمة} (Abstractions) فوق العتاد المادي:
\begin{itemize}
  \item بدلاً من السماح للبرامج بالوصول المباشر للذاكرة الفيزيائية، يوفر نظام التشغيل لكل عملية \textbf{فضاء العناوين الافتراضي} (Virtual Address Space) خاصًا بها.
  \item بدلاً من الوصول المباشر لقطاعات القرص الصلب، يقدم نظام التشغيل تجريد \textbf{نظام الملفات} (File System).
  \item بدلاً من التحكم المباشر بأجهزة الإدخال والإخراج، توفر النواة واجهة واصفات الملفات.
  \item بدلاً من التحكم المباشر بأنوية المعالج، يوفر نظام التشغيل تجريد \textbf{العملية} (Process) والجدولة الزمنية.
\end{itemize}

هذه التجريدات ليست فقط تجعل كتابة البرامج أسهل، بل تضمن العزل المحكم؛
لأن البرامج تتفاعل فقط مع التجريدات النواتية وليس مع الموارد المادية المباشرة،
مما يتيح للنواة تفحص وتفتيش كل طلب ومنع أي تجاوز للحدود المسموح بها.

\section{وضع المستخدم، وضع المشرف، واستدعاءات النظام}
\label{sec:user_supervisor_syscalls}

يتطلب تحقيق العزل القوي دعمًا صريحًا ومحددًا من عتاد وحدة المعالجة المركزية (CPU).
توفر معمارية \textsf{RISC-V} ثلاثة مستويات من الصلاحيات والأنماط التي ينفذ فيها المعالج التعليمات:
\begin{enumerate}
  \item \textbf{وضع الآلة} (Machine Mode - M-mode): أعلى مستوى صلاحية، يمتلك وصولاً كاملاً وغير مشروط للعتاد المادي، ويُستخدم عادة لتهيئة العتاد عند الإقلاع الأول (مثل البرمجيات الثابتة Firmware / BIOS).
  \item \textbf{وضع المشرف} (Supervisor Mode - S-mode): المستوى المخصص لنواة نظام التشغيل. يسمح بتنفيذ التعليمات الممتازة (Privileged Instructions) مثل تعديل سجلات إدارة الذاكرة وتفعيل المقاطعات وتحويل وضعيات المعالج.
  \item \textbf{وضع المستخدم} (User Mode - U-mode): المستوى المخصص لبرامج المستخدم العادية. يُمنع المعالج في هذا الوضع من تنفيذ التعليمات الممتازة أو الوصول المباشر للسجلات العتادية الحساسة.
\end{enumerate}

عندما تحاول عملية تعمل في وضع المستخدم تنفيذ تعليمة ممتازة (مثل تعديل سجل التحكم بالتصفح \texttt{satp})، يرفض العتاد التعليمة ولا ينفذها ويولّد \textbf{استثناء} (Exception) يُحاط به فورًا وتنتقل السيطرة إلى معالج النواة.

للإنتقال المنظم من وضع المستخدم إلى وضع المشرف، يوفر عتاد RISC-V تعليمة خاصة تُدعى \texttt{ecall} (Environment Call).
عند تنفيذ \texttt{ecall}:
\begin{enumerate}
  \item يحفظ المعالج حالة البرنامج الحالية ويعيد رفع مستوى الصلاحية إلى وضع المشرف (S-mode).
  \item يقفز المعالج إلى عنوان معالج المصايد (Trap Handler) المسجل مسبقًا في السجل الممتاز \texttt{stvec}.
  \item تنفذ النواة الخدمة المطلوبة وتتحقق من صحة المدخلات.
  \item تعود النواة إلى وضع المستخدم باستخدام تعليمة \texttt{sret} (Supervisor Return).
\end{enumerate}

\section{تنظيم النواة}
\label{sec:kernel_org}

السؤال الهيكلي في تصميم نظام التشغيل هو: ما الأجزاء التي ينبغي أن تعمل في وضع المشرف (S-mode) وتكون داخل النواة؟
تتوزع أنظمة التشغيل في هيكلة النواة إلى نمطين رئيسين:
\begin{itemize}
  \item \textbf{النواة المتجانسة} (Monolithic Kernel): تنفذ جميع مكونات نظام التشغيل (إدارة الذاكرة، الجدولة، برامج تشغيل الأجهزة، نظام الملفات، الأنابيب، شبكة الاتصالات) داخل فضاء النواة وبصلاحيات كاملة (S-mode). يوفر هذا التصميم أداءً عاليًا وسرعة في استدعاءات الخدمات لأن جميع الخدمات تعمل في فضاء واحد ولا تتطلب تبديل السياق. يعتمد نظام \texttt{xv6}، مثل Linux و FreeBSD، هيكلة النواة المتجانسة.
  \item \textbf{نواة مصغّرة} (Microkernel): تنفذ النواة الحد الأدنى من الخدمات (الجدولة، الاتصال بين العمليات IPC، إدارة الذاكرة المباشرة)، بينما تنفذ بقية الخدمات مثل نظام الملفات وبرامج تشغيل الأجهزة كعمليات عادية في فضاء المستخدم (U-mode). يزيد هذا التصميم من أمان النظام واستقراره لأن خلل برامج تشغيل الأجهزة لا ينهي النظام بأكمله، ولكنه يتطلب تكلفة إضافية في تبديل السياق والاتصالات بين العمليات.
\end{itemize}

\section{الكود البرمجي: هيكلة xv6}
\label{sec:code_xv6_org}

تتواجد الشفرة المصدرية لنواة \texttt{xv6} داخل الدليل \texttt{kernel/}.
وتُقسم إلى وحدات برمجية واضحة:
\begin{itemize}
  \item \texttt{kernel/main.c}: نقطة بداية تشغيل النواة عند الإقلاع وتوزيع المهام الأولية.
  \item \texttt{kernel/proc.c}: إدارة العمليات، والجدولة، وإنشاء وتدمير العمليات، وتغيير السياق.
  \item \texttt{kernel/vm.c}: إدارة الذاكرة الافتراضية وجداول الصفحات وتعيين العناوين.
  \item \texttt{kernel/kalloc.c}: مخصص الذاكرة المادية والصفحات الحرة.
  \item \texttt{kernel/trap.c}: معالجة المصايد، الاستثناءات، المقاطعات العتادية، واستدعاءات النظام.
  \item \texttt{kernel/spinlock.c} و \texttt{kernel/sleeplock.c}: أقفال التناوب وأقفال النوم للتزامن وحماية هياكل البيانات.
  \item \texttt{kernel/bio.c} و \texttt{kernel/fs.c} و \texttt{kernel/log.c}: طبقات المخزن المؤقت ونظام الملفات وسجل المعاملات.
  \item \texttt{kernel/file.c} و \texttt{kernel/pipe.c} و \texttt{kernel/sysfile.c}: إدارة واصفات الملفات والأنابيب واستدعاءات نظام الملفات.
  \item \texttt{kernel/uart.c} و \texttt{kernel/virtio\_disk.c}: برامج تشغيل الأجهزة للمنفس التسلسلي UART والقرص المادي.
\end{itemize}

\section{نظرة عامة على العمليات}
\label{sec:process_overview}

تُمثَّل كل عملية في \texttt{xv6} بواسطة هيكل بيانات C يُدعى \texttt{struct proc} (المعرف في \texttt{kernel/proc.h}).
يحتفظ هذا الهيكل بجميع معلومات حالة العملية:
\begin{itemize}
  \item حالة العملية (\texttt{enum procstate}): هل هي \texttt{UNUSED}، \texttt{EMBRYO}، \texttt{SLEEPING}، \texttt{RUNNABLE}، \texttt{RUNNING}، أو \texttt{ZOMBIE}.
  \item معرّف العملية (\texttt{pid}).
  \item جدول الصفحات الخاص بالعملية (\texttt{pagetable\_t pagetable}).
  \item المكدس النواتي الخاص بالعملية (\texttt{uint64 kstack}).
  \item إطار المصيدة (\texttt{struct trapframe *trapframe}) للحفظ والاسترجاع لسجلات المستخدم.
  \item سياق العملية (\texttt{struct context context}) لحفظ سجلات النواة أثناء تبديل المجدول.
  \item قائمة الملفات المفتوحة (\texttt{struct file *ofile[NOFILE]}).
  \item الدليل الحالي للعملية (\texttt{struct inode *cwd}).
\end{itemize}

لكل عملية مكدسان اثنان: مكدس المستخدم (User Stack) ومكدس النواة (Kernel Stack).
عندما تنفذ العملية تعليمات في فضاء المستخدم، يُستخدم مكدس المستخدم.
وعندما تتنقل العملية إلى النواة عبر مصيدة أو استدعاء نظام، تحول النواة التنفيذ فورًا إلى مكدس النواة الخاص بهذه العملية والمخصص في ذاكرة المشرف.

\section{الكود البرمجي: إقلاع xv6، العملية الأولى واستدعاء النظام}
\label{sec:code_boot_first_proc}

عند تشغيل الحاسوب (أو محاكي QEMU):
\begin{enumerate}
  \item يبدأ المعالج في تنفيذ كود الـ ROM الإقلاعي في وضع الآلة (M-mode)، والذي يحمل النواة إلى الذاكرة عند العنوان الفيزيائي \texttt{0x80000000}.
  \item ينفذ المعالج كود التجميع في \texttt{kernel/entry.S}، لتهيئ مكدس النواة الابتدائي لكل نوات معالجة (Hart).
  \item تنتقل السيطرة إلى الدالة \texttt{main()} في \texttt{kernel/main.c} في وضع المشرف (S-mode).
  \item تقوم النواة بتهيئة الأقفال، ووحدة إدارة الذاكرة، ومجمع المقاطعات PLIC، والذاكرة المخبئية.
  \item تدعو \texttt{main()} الدالة \texttt{userinit()} لإنشاء أول عملية مستخدم.
  \item تنفذ العملية الأولى كوداً مجمعاً صغيراً بداخل \texttt{initcode.S} يتضمن استدعاء النظام \texttt{exec("/init")} ليعمل برنامج البدء الرسمي الشارح لمغلف الأوامر \texttt{sh}.
\end{enumerate}

\section{نموذج الأمان}
\label{sec:security_model}

يفترض نموذج الأمان في \texttt{xv6} أن النواة موثوقة تمامًا ومحصنة ضد الثغرات البرمجية، بينما برامج المستخدم غير موثوقة ومفصولة تمامًا. تقوم النواة بالتحقق من جميع العناوين الممررة من برامج المستخدم للتأكد من وقوعها ضمن النطاق المسموح به للعملية قبل إجراء القراءة أو الكتابة.

\section{العالم الحقيقي}
\label{sec:real_world_first}

تستخدم أنظمة التشغيل الإنتاجية مثل Linux آليات أمان متطورة مثل مساحات الأسماء (Namespaces) والحاويات (Containers) وASLR (عشوائية تخطيط فضاء العناوين) لتعزيز العزل. ومع ذلك، تشترك جميع هذه الأنظمة مع \texttt{xv6} في الأعمدة الثنائية ذاتها: العزل العتادي بين وضع المستخدم ووضع المشرف، واستخدام استدعاءات النظام المحكومة عبر التعليمة \texttt{ecall}.

\section{تمارين}
\label{sec:exercises_first}

\begin{enumerate}
  \item تتبع خطوات تنفيذ تعليمة \texttt{ecall} بدءًا من وضع المستخدم وصولاً إلى معالج \texttt{syscall()} في \texttt{kernel/syscall.c}.
  \item أضف استدعاء نظام جديد باسم \texttt{trace(mask)} يقوم بطباعة أسماء استدعاءات النظام المنفذة لاحقاً.
\end{enumerate}
